\subsection{Conditional mean, variance, and quantile}\label{sec:mean}
The protocol uses three numerical features of the truncated distribution. The conditional mean and the conditional variance
describe the direction and the dispersion of the shift of the estimate, and the lower quantile directly determines the threshold of
the second-round capacity. This subsection derives their expressions one by one from the limiting distribution of
Theorem~\ref{thm:main}.

\paragraph{The first two moments.}
Theorem~\ref{thm:main} gives uniform convergence of the distribution functions. Lifting this to convergence of moments requires
uniform integrability. For every $M$, $\E X^2=1$ holds exactly, and the limiting failure probability $\Phi(z_r)$ is positive; from
this one verifies the uniform integrability of $X^2$ under failure (Appendix~\ref{app:normalization}), and the first two conditional
moments converge to the corresponding values of the limiting distribution.

The limiting distribution is the distribution of the first coordinate of a bivariate normal truncated to the half-plane
$\{y\le z_r\}$; its moments can be expressed through the truncated moments of a standard normal. Let $A(z)=\phi(z)/\Phi(z)$ denote
the inverse Mills ratio. The mean of a standard normal variable truncated to $\{y\le z\}$ is $-A(z)$, and its variance is
$1-A(z)[z+A(z)]$. In the limiting bivariate normal, the first coordinate decomposes into $c$ times the standardized score plus a
normal term of variance $1-c^2$ that is independent of the score. The truncation acts only on the score, leaving the independent
term unaffected. Hence
\begin{equation*}
\E[X\mid F]\longrightarrow-cA(z_r),\qquad
\Var(X\mid F)\longrightarrow1-c^2A(z_r)\bigl[z_r+A(z_r)\bigr].
\end{equation*}
The shift of the mean equals the correlation coefficient $c$ times the average downward shift $A(z_r)$ of the score after
truncation: the part of the measurement proportional to the score moves together with the score, rightward when $c<0$ and leftward
when $c>0$. The conditional variance is less than one, because truncation compresses the fluctuations along the score direction
while the part independent of the score is preserved unchanged. Under balanced signs $c=0$, both return to their unconditional
values, consistent with Corollary~\ref{cor:balanced}. The smaller $z_r$ is (that is, the higher the load), the lower the failure
probability, the larger $A(z_r)$, and the more pronounced the shift.

\paragraph{Quantile.}
Let $q_{\rm sel}(u;z,c)$ denote the quantile at level $u\in(0,1)$ of the distribution function $\Phi_2(x,z;c)/\Phi(z)$. This
distribution has density
\begin{equation*}
f_{\rm sel}(x;z,c)=\frac{\phi(x)}{\Phi(z)}
\Phi\!\left(\frac{z-cx}{\sqrt{1-c^2}}\right)
\end{equation*}
which is positive everywhere, and the distribution function is strictly increasing, so the conditional quantile at every fixed level
converges to $q_{\rm sel}$.

Which side the quantile moves toward is likewise determined by the sign of $c$. Since
$\partial_c\Phi_2(x,z;c)=\phi_2(x,z;c)>0$, at every fixed $x$ the conditional distribution function increases strictly in $c$; at
$c=0$ it is the standard normal distribution function. Lifting the distribution function overall means that the quantile at every
level shifts leftward. Thus when $c>0$ all limiting quantiles lie below the corresponding standard normal values, and when $c<0$
they all lie above.

Thus far the failure-conditioned law inside the critical window is determined by only two numbers, $z_r$ and
$c=\eta_k\rho_\infty^2$. The direction of the shift of the mean, the variance, and the quantile all depend on the sign of $\eta_k$.
